% Satu berkas per bab. \chapter menghasilkan "BAB I" lalu judulnya di baris berikut.
\chapter{Pendahuluan}\label{bab:pendahuluan}

\section{Latar Belakang}\label{sec:latar-belakang}

Paragraf pertama setiap bab dan subbab tidak perlu perlakuan khusus; indentasi 0,75~cm
sudah diatur gaya. Rujukan silang ditulis \texttt{Bab~\textbackslash ref\{bab:pendahuluan\}}
sehingga menghasilkan Bab~\ref{bab:pendahuluan}.

Sitasi memakai biblatex gaya \emph{authoryear}: \parencite{contoh2020} untuk sitasi dalam
kurung, dan \textcite{contoh2020} bila nama penulis menjadi bagian kalimat.

Istilah asing dicetak miring, misalnya \textit{deep learning}. Tanda yang masih harus
dilengkapi ditandai agar mudah dicari: \todo{lengkapi angka ini},
\rujukandiperlukan{klaim ini}, dan \konfirmasi{tanyakan ke pembimbing}.

\section{Rumusan Masalah}\label{sec:rumusan-masalah}

\begin{enumerate}[label=RM\arabic*.,leftmargin=1.2cm,labelwidth=1.2cm]
  \item Pertanyaan penelitian pertama.
  \item Pertanyaan penelitian kedua.
\end{enumerate}

\section{Contoh Tabel}\label{sec:contoh-tabel}

Keterangan tabel diletakkan di atas tabel, keterangan gambar di bawah gambar.

\begin{table}[H]
  \centering
  \caption{Keterangan tabel ditulis di atas}
  \label{tab:contoh}
  \begin{tabular}{|l|l|}
    \hline
    \textbf{Kolom pertama} & \textbf{Kolom kedua} \\ \hline
    Isi baris pertama      & Isi baris pertama    \\ \hline
    Isi baris kedua        & Isi baris kedua      \\ \hline
  \end{tabular}
  \sumber{Penulis, \tahun.}
\end{table}

Untuk tabel yang melampaui satu halaman, pakai \texttt{xltabular} dan perintah
\texttt{\textbackslash lanjutan\{n\}} pada baris penanda sambungan.

\section{Contoh Persamaan}\label{sec:contoh-persamaan}

Persamaan bernomor otomatis mengikuti nomor bab.

\begin{equation}\label{eq:contoh}
  y = \sum_{i=1}^{n} w_i x_i + b
\end{equation}

Persamaan~\ref{eq:contoh} dirujuk seperti ini.
